\subsection{Equivalent Interest Rates} 
One reason we might do this is to think of interest as compounded continuously.%, which is easier for computation, or annually, which is easier to understand intuitively.
\begin{example}
Suppose we take out a loan with 6.5\% interest, compounded quarterly. What would be an equivalent interest rate compounded continuously? Give your answer as a percent, rounded to two decimal places.
\begin{lpic}{}{4}
\end{lpic}
The given interest rate is equivalent to \makeblanksmall\% interest compounded continuously.
\end{example}
We can also work in the other direction on compound interest, solving for $r$.
\begin{example}
Suppose we take out a loan with 8.2\% interest, compounded continuously. What would be an equivalent interest rate compounded annually? Give your answer as a percent, rounded to two decimal places.
\begin{lpic}{}{3}
\end{lpic}
The given interest rate is equivalent to \makeblanksmall\% interest compounded annually.
\end{example}